% ATTEMPT_STATUS: unresolved
% RESEARCH_GENERATED: 2026-10-01; GPT-6 Astra Ultra; individual mathematical attempt
% TOP_PROBLEM: {"id": 126, "title": "Difference Sets Containing Squares", "category_id": 1, "category": {"id": 1, "name": "number_theory", "display_name": "Number Theory", "description": "Properties of integers, prime numbers, Diophantine equations.", "slug": "number-theory", "order_index": 1, "created_at": "2026-07-31T15:26:25.670Z"}, "status": "open"}
% FIRST_POSTED: 2026-10-01
% Imported from: unsolvedmath_additional_500.tex; UnsolvedMath ID 126
% !TeX program = lualatex
% Compile twice: lualatex 126.tex
% Self-contained: no external bibliography, images, or \input files.
% Numeric section labels and visible numbers are original UnsolvedMath IDs.
\documentclass[11pt,a4paper]{article}
\usepackage[margin=24mm,headheight=14pt]{geometry}
\usepackage{fontspec}
\setmainfont{Latin Modern Roman}
\setsansfont{Latin Modern Sans}
\setmonofont{Latin Modern Mono}
\usepackage{amsmath,amssymb,mathtools}
\usepackage{unicode-math}
\setmathfont{Latin Modern Math}
\usepackage{microtype}
\usepackage{enumitem,xurl,needspace}
\usepackage[dvipsnames]{xcolor}
\usepackage{hyperref}
\hypersetup{unicode=true,colorlinks=true,linkcolor=MidnightBlue,urlcolor=MidnightBlue,
 pdftitle={UnsolvedMath Problem 126},pdfauthor={Source-annotated compilation},bookmarksnumbered=true}
\usepackage{bookmark,tocloft,titlesec,fancyhdr}
\setlength{\cftsecnumwidth}{4.2em}
\cftsetpnumwidth{2.5em}
\cftsetrmarg{4em}
\titleformat{\section}{\Large\bfseries\raggedright}{\thesection}{0.7em}{}
\pagestyle{fancy}\fancyhf{}
\fancyhead[L]{\small\sffamily UnsolvedMath research attempt}
\fancyhead[R]{\small Original catalogue IDs}
\fancyfoot[C]{\thepage}
\setlength{\parindent}{0pt}
\setlength{\parskip}{5pt plus 1pt}
\setlength{\emergencystretch}{3em}
\setcounter{tocdepth}{1}\setcounter{secnumdepth}{1}
\setlist[itemize]{leftmargin=3em,itemsep=4pt,topsep=4pt,parsep=0pt}
\allowdisplaybreaks
\newcommand{\N}{\mathbb{N}}
\newcommand{\Z}{\mathbb{Z}}
\newcommand{\Q}{\mathbb{Q}}
\newcommand{\R}{\mathbb{R}}
\newcommand{\C}{\mathbb{C}}
\newcommand{\F}{\mathbb{F}}
\newcommand{\A}{\mathbb{A}}
\newcommand{\PP}{\mathbb{P}}
\newcommand{\E}{\mathbb{E}}
\newcommand{\eps}{\varepsilon}
\newcommand{\dd}{\,\mathrm d}
\newcommand{\sref}[2]{\hyperlink{source:#1:#2}{[S#2]}}
\newif\ifshowcatalogue
\showcataloguetrue
\newcommand{\cataloguescope}[1]{\ifshowcatalogue
 \paragraph{Statement recorded in the supplied source.}
 #1\fi}
\begin{document}
% UnsolvedMath ID 126; source catalogue code GREEN-038

\setcounter{section}{125}

\section{A Square in a Dense Difference Set}\label{126}

{\small\sffamily UnsolvedMath ID 126 · Number Theory}

\subsection{Definitions and mathematical statement}

\paragraph{Shared notation.}Unless a section says otherwise, $\N=\{1,2,\ldots\}$,
$\Z$ is the ring of integers, $\Q,\R,\C$ are the rational, real, and complex
fields, $[N]=\{1,\ldots,\lfloor N\rfloor\}$, and $\log$ is the natural
logarithm. For real functions of a parameter tending to infinity,
$f\ll g$ and $f=O(g)$ mean $|f|\leq Cg$ eventually for some fixed $C$;
$f\gg g$ reverses this comparison; $f=o(g)$ means $f/g\to0$;
$f\sim g$ means $f/g\to1$; and $f\asymp g$ means both order bounds.
A subscript on $O$ or $\ll$ specifies allowed constant dependence.
The expression $n^{o(1)}$ in an upper bound means growth smaller than
$n^\epsilon$ for each fixed $\epsilon>0$. Local definitions govern any
exception, finite-field convention, or use of the same symbol for another
object. An asymptotic claim is not automatically an effective algorithm.



For a finite set $A\subseteq[N]$, $A-A=\{a-b:a,b\in A\}$. A nonzero square is $u^2$ with $u\in\N$. The asymptotic assertion seeks an absolute $c>0$ so that sets at the indicated power-size threshold contain $a,b$ with $a-b=u^2>0$. Cardinality thresholds are rounded to integers; the chosen $c$ must not depend on $N$. \sref{126}{1} \sref{126}{2}

\paragraph{Statement recorded in the supplied source.}
Is there $c > 0$ such that whenever $A \subset [N]$ has size $N^{1-c}$, the difference set $A - A$ contains a nonzero square? \sref{126}{1}

\paragraph{Residual task reported by the archive.}

The embedded review (2026-08-17) records: Prove any fixed polynomial saving for square-difference-free subsets of [N], or construct examples ruling it out. \sref{126}{1}

\subsection{Short English statement}

Does a fixed power-saving density threshold already force two selected integers to differ by a nonzero perfect square?

\subsection{Research attempt}

\paragraph{Provenance and scope.}
This individual attempt was generated by GPT-6 Astra Ultra on 1 October 2026. The method was to derive explicit partial results and test the first proposed extension adversarially. The original statement and sources below are retained. No claim of mathematical novelty or complete solution is made.

\paragraph{Formal target and primary-source context.}
A square-difference-free set $A\subset[N]$ has no two elements with positive difference $u^2$, $u\geq1$. The question asks for a fixed power saving in its maximum cardinality. The inspected Green--Sawhney manuscript proves a bound of shape $Ne^{-c\sqrt{\log N}}$. That remains larger than $N^{1-\delta}$ for every fixed $\delta>0$ at sufficiently large $N$, since $\sqrt{\log N}=o(\log N)$. We construct a simple digit family and derive an exact finite-field upper bound, then check why reducing modulo a prime fails to connect the two at the requested scale.

\paragraph{A square-free difference construction in base five.}
For $k\geq1$, let
\[
 A_k=\left\{\sum_{i=0}^{k-1}\varepsilon_i5^i:\varepsilon_i\in\{0,2\}\right\}.
\]
It has $2^k$ elements and maximum $(5^k-1)/2$. Take distinct $a,b\in A_k$, and let $j$ be their lowest differing digit. Then
\[
 a-b=5^j u,\qquad u\equiv2\text{ or }-2\pmod5.
\]
There is no ambiguity from higher digits: after factoring out $5^j$, they are multiples of five, while the first nonzero digit difference is $\pm2$. If $a-b$ were a nonzero integer square, its five-adic valuation $j$ would be even and its unit part $u$ would be a nonzero square modulo five. But $2$ and $3=-2$ are both nonsquares modulo five. This contradiction proves that $A_k$ has no nonzero square difference of either sign. Translating by one puts it in a positive interval without changing differences.

This yields examples of cardinality order $N^{\log_5 2}$ along an explicit sequence of lengths. Therefore a proposed universal power exponent cannot exceed $1-\log_5 2$ by this construction alone. This is not asserted to be the strongest known construction; it is a transparent example with a complete valuation proof.

\paragraph{An exact quadratic-character bound over prime fields.}
Let $p\equiv1\pmod4$ be prime and $B\subset\F_p$ have all nonzero differences quadratic nonresidues. Let $m=|B|$ and let $\chi$ be the quadratic character, with $\chi(0)=0$. Then
\[
 \sum_{a,b\in B}\chi(a-b)=-m(m-1).
\]
For a nonzero additive frequency $r$, the Fourier transform of $\chi$ has modulus $\sqrt p$, by the usual Gauss-sum identity; its zero coefficient is zero. A short proof of the relevant modulus follows from $\sum_x\chi(x)\chi(x+h)=-1$ for $h\ne0$. To check this correlation, divide by $h$, complete the square, and count solutions to $v^2=u(u+1)$: multiplying by four converts it to $(2u+1)^2-(2v)^2=1$, whose factor pairs number $p-1$. Thus the correlation is $-1$, and Fourier inversion gives modulus $\sqrt p$ at every nonzero frequency.

The convolution operator with kernel $\chi$ therefore has operator norm $\sqrt p$ on mean-zero functions. Apply it to $1_B-m/p$; its squared norm is $m-m^2/p$ in counting measure. We obtain
\[
 m(m-1)\leq\sqrt p\,(m-m^2/p),
 \qquad m\leq\sqrt p.
\]
For $m>0$, the final inequality follows by division and rearrangement; the empty set is immediate. This is a strong power bound in a finite-field model.

\paragraph{The attempted modular transfer breaks.}
An integer square-difference-free set need not have only quadratic-nonresidue differences modulo $p$. A positive integer can be a square modulo $p$ without being an integer square. For example, two integers differing by six have no square difference, but six is a square modulo five. Thus the finite-field character condition is strictly stronger than the original one. Choosing a prime comparable to $N$ does not repair this implication. A modular proof would need extra information forcing an actual integer square among the differences, not merely a quadratic-residue class.

\paragraph{Verification and remaining gap.}
The companion script checks every pair in the digit construction for $k=1,\ldots,5$, obtaining sizes $2,4,8,16,32$ and no square differences. The prime-field argument explicitly uses $p\equiv1\pmod4$, ensuring the undirected difference condition is consistent. Its normalisations and diagonal zero terms are retained. The work proves an explicit infinite lower family, a finite-field power bound, and the failure of the direct reduction between them. The unresolved task is an integer upper bound with any fixed exponent saving; neither the current subexponential literature bound nor the stronger modular hypothesis supplies it.

\subsection{Sources}

{\small

\begin{itemize}

\item[\hypertarget{source:126:1}{S1}] ULAM.AI, \emph{UnsolvedMath}, supplied \texttt{problems.json}, record 126 (GREEN-038), statement and background. The embedded literature-review date is 2026-08-17; that is the archive’s review date, not a fresh certification. Dataset landing page: \url{https://huggingface.co/datasets/ulamai/UnsolvedMath}.

\item[\hypertarget{source:126:2}{S2}] Ben Green and Mehtaab Sawhney, New bounds for the Furstenberg-Sarkozy theorem, arXiv:2411.17448 (2025 version). \url{https://arxiv.org/abs/2411.17448}. Bibliographic information imported from the supplied record.

\item[\hypertarget{source:126:3}{S3}] Ben Green, 100 Open Problems, current Problem 65. \url{https://people.maths.ox.ac.uk/greenbj/papers/open-problems.pdf}. The author-hosted document was inspected on 27 September 2026; the archive’s local code is not a problem-number locator in this paper.

\item[E1] B. Green and M. Sawhney, \emph{New bounds for the Furstenberg--Sarkozy theorem}, current manuscript. \url{https://arxiv.org/abs/2411.17448}. Inspected 1 October 2026.

\end{itemize}

}

\label{end:126}
\end{document}
